\section*{Chapitre \ref{ch:dv:build-a-pricing-library} --- Construction~: une bibliothèque de valorisation}

\begin{solution}{exo:dv:build-a-pricing-library:1}
Python introduit un sel aléatoire dans le \texttt{hash()} des chaînes à chaque processus, si bien que le même instrument recevrait des graines différentes dans deux exécutions ou deux processus de calcul, et un bump
calculé entre processus comparerait des nombres aléatoires différents. Une somme de contrôle (CRC-32) de l'identifiant est la même partout et pour toujours.
\end{solution}

\begin{solution}{exo:dv:build-a-pricing-library:2}
Chacune est valorisée à la volatilité de son propre strike, 100 à un an, qui est un nœud. Le bump du nœud 90 ne modifie la surface
qu'entre 80 et 100 (interpolation linéaire en log-strike), et non en 100 lui-même. La barrière à 80 du knock-out n'aurait d'effet que pour un moteur qui lit les volatilités le long de la trajectoire.
\end{solution}

\begin{solution}{exo:dv:build-a-pricing-library:3}
Parité in-out (les deux barrières ensemble valent l'option vanille)~; le call et le put digitaux ensemble versent une unité de paiement certaine,
valant le facteur d'actualisation~; sans effet de levier et avec un objectif inatteignable, le TARF est un strip de puts achetés, un par date de fixing.
\end{solution}

\begin{solution}{exo:dv:build-a-pricing-library:4}
Avec des nombres aléatoires communs, les calculs avec et sans bump ne diffèrent que par le bump~: il reste donc l'erreur de Monte Carlo de l'estimateur
de sensibilité lui-même, qui décroît comme $1/\sqrt n$. Pour la réduire, augmenter le nombre de trajectoires, utiliser une variable de contrôle, un payoff plus régulier près des barrières, ou un \omterm{def:dv:greeks-and-the-hedging-pnl:greeks}{vega} pathwise ou par rapport de vraisemblance (chapitre 23)
à partir des \texttt{greeks} propres au moteur.
\end{solution}

\begin{solution}{exo:dv:build-a-pricing-library:5}
Entre les échéances, la surface interpole la variance totale, $\sigma^2t$, et non la volatilité~; à échéance donnée, elle interpole en log-strike. Un décalage
parallèle de toutes les volatilités n'est pas exactement la somme des décalages de ses nœuds. Les différences centrées laissent en outre des termes du second ordre, qui diffèrent selon qu'on applique un bump
large ou plusieurs petits.
\end{solution}

\begin{solution}{exo:dv:build-a-pricing-library:6}
Le prix du moteur ne suit plus l'instantané~: les bumps et les scénarios ne l'atteignent plus, ses \omterm{def:dv:greeks-and-the-hedging-pnl:greeks}{grecques} sortent nulles, et deux threads valorisant des instantanés différents
interfèrent entre eux. Les lots parallèles deviennent faux sans qu'aucune erreur ne soit levée.
\end{solution}

\begin{solution}{exo:dv:build-a-pricing-library:7}
Avec 400\,000 trajectoires, l'écart type sur vingt graines vaut 25.1 (moyenne 2\,231), contre 50.8 avec 100\,000~: la moitié, comme le prédit $1/\sqrt n$ pour quatre fois plus de trajectoires.
\end{solution}

\begin{solution}{exo:dv:build-a-pricing-library:8}
L'accord avec l'ancien moteur montre seulement que les deux donnent les mêmes résultats, et non que l'un ou l'autre a raison~: ils peuvent partager un bug, une convention ou une donnée de marché. Le test doit
être confronté à une référence indépendante (une formule fermée, une grille convergée) et à une identité, sur toute la plage des paramètres, y compris près des barrières, des échéances
et profondément dans la monnaie comme hors de la monnaie.
\end{solution}

\begin{solution}{pb:dv:build-a-pricing-library:1}
\textbf{1.} Call~: Analytic~; put américain~: Tree~; knock-out~: ClosedForm (surveillance quotidienne par le décalage de Broadie--Glasserman--Kou)~; \omterm{def:dv:variance-swaps-and-volatility-derivatives:vs}{swap de variance}~: ClosedForm
(réplication sur la surface)~; \omterm{def:dv:autocallables:autocallable}{autocallable}~: MonteCarlo. Le registre renvoie le premier moteur enregistré qui prend en charge la paire.
\textbf{2.} $-928\,523$.
\textbf{3.} Analytique et arbre à $2\times10^{-6}$ près~; la grille à 0.0024 (200 nœuds)~; Monte Carlo dans son erreur type.
\textbf{4.} Lire la volatilité sur la surface de l'instantané, aux strikes et aux dates qui lui sont propres, jamais à partir d'un nombre stocké.
\textbf{5.} Il lève une erreur de valorisation qui nomme l'instrument et le modèle.
\textbf{6.} Call 0, put américain $-88$, knock-out 0, \omterm{def:dv:variance-swaps-and-volatility-derivatives:vs}{swap de variance} $-90$, \omterm{def:dv:autocallables:autocallable}{autocallable} $+2\,249$~: 2\,072.
\textbf{7.} Son moteur lit la variance totale à 90, au milieu entre la barrière de knock-in et le déclencheur, et le titre vaut moins pour son porteur lorsque la volatilité
monte~; la banque est vendeuse du titre, donc longue de son \omterm{def:dv:greeks-and-the-hedging-pnl:greeks}{vega}.
\textbf{8.} L'interpolation linéaire en log-strike entre 90 et 100 place environ la moitié du poids d'un strike de 95 sur le nœud 90.
\textbf{9.} $-90$~: le strip d'options qui réplique la variance couvre tous les strikes, si bien que chaque nœud porte une partie de son \omterm{def:dv:greeks-and-the-hedging-pnl:greeks}{vega}.
\textbf{10.} 2\,275 à un an, 1\,464 à deux ans et 3\,008 à trois ans par point, et rien avant un an.
\textbf{11.} 51 avec des nombres aléatoires communs, 217 sans.
\textbf{12.} Les prix avec et sans bump partagent leurs erreurs, qui s'annulent dans la différence~; sans partage, les erreurs de deux calculs indépendants
s'additionnent et sont divisées par le bump.
\textbf{13.} Environ la moitié~: 25.
\textbf{14.} Pour que le prix d'une transaction ne change pas lorsque d'autres transactions sont ajoutées au portefeuille ou en sont retirées, et pour que le livre~6 puisse revaloriser n'importe quel sous-ensemble.
\textbf{15.} Le portefeuille est long de volatilité et long d'un krach ($+161.8$ milliers pour $-20\%$ avec dix points de volatilité en plus), et perd sur les hausses
($-24.0$ milliers à $+10\%$) et sur la baisse de la volatilité ($-34.9$ milliers).
\textbf{16.} À une centaine près~: la seule partie Monte Carlo a un écart type de 51.
\textbf{17.} Les prix de référence et les contrôles croisés entre moteurs sont exécutés chaque nuit, avec un historique du moteur et de la graine utilisés pour chaque
transaction, l'instantané stocké avec chaque résultat, et des limites sur les variations non expliquées.
\textbf{18.} Un test de propriété qui applique un bump à l'instantané et exige que le prix change pour chaque instrument dont la valeur dépend du facteur choqué.
\textbf{19.} 2\,072 par point de volatilité (\omterm{def:dv:autocallables:autocallable}{autocallable} $+2\,249$, put américain $-88$, \omterm{def:dv:variance-swaps-and-volatility-derivatives:vs}{swap de variance} $-90$), avec un écart type de 51 dû au bruit de Monte Carlo
avec des nombres aléatoires communs, et 217 sans.
\textbf{20.} Valoriser comme une fonction pure d'un instantané immuable, afin qu'un instantané soumis à un bump revalorise tout, avec des nombres aléatoires partagés pour la simulation.
\end{solution}

\begin{solution}{iq:dv:build-a-pricing-library:1}
Pour que chacun varie indépendamment~: un nouveau produit n'a besoin que de ses termes et d'un moteur~; une nouvelle méthode valorise chaque produit qu'elle prend en charge~;
deux moteurs se contrôlent mutuellement~; les hypothèses du modèle sont explicites et interchangeables.
\iqlookfor{variation indépendante, contrôle croisé.}
\end{solution}

\begin{solution}{iq:dv:build-a-pricing-library:2}
Nombres aléatoires communs~: la même graine et la même grille pour chaque revalorisation d'un instrument, si bien que les bumps éliminent par différence l'erreur de simulation.
Suivent ensuite des payoffs plus réguliers et des \omterm{def:dv:greeks-and-the-hedging-pnl:greeks}{grecques} directes.
\iqlookfor{la même graine par instrument.}
\end{solution}

\begin{solution}{iq:dv:build-a-pricing-library:3}
Des instantanés immuables pour le risque~: la valorisation est une fonction pure, les lots parallèles sont sûrs, les scénarios sont de nouveaux instantanés, et les résultats sont reproductibles à partir des entrées stockées.
Les observateurs et la mise en cache conviennent à un écran de trading qui revalorise à chaque variation de cotation~; ils compliquent le calcul sur des milliers de marchés choqués et l'évaluation parallèle.
\iqlookfor{pureté et reproductibilité face au recalcul incrémental.}
\end{solution}

\begin{solution}{iq:dv:build-a-pricing-library:4}
Par confrontation à des formules fermées et à des \omterm{def:dv:build-a-pricing-library:reference}{pricers de référence} indépendants, à d'autres moteurs, par des identités (parité, in-out), par des cas limites (volatilité nulle, strikes lointains), par convergence
selon ses propres paramètres (pas, trajectoires), et par des tests de régression ancrés sur des valeurs certifiées.
\iqlookfor{références indépendantes et convergence.}
\end{solution}

\begin{solution}{iq:dv:build-a-pricing-library:5}
Placer chaque moteur sur le même objet surface, avec des facteurs de risque par nœud~; appliquer un bump à chaque nœud à tour de rôle sur une copie de l'instantané, revaloriser chaque position
avec son propre modèle et son propre moteur, avec des nombres aléatoires communs pour les simulations, puis sommer par nœud. Vérifier que la somme des nœuds redonne le \omterm{def:dv:greeks-and-the-hedging-pnl:greeks}{vega} parallèle.
\iqlookfor{une seule surface, bumps par nœud, nombres aléatoires communs.}
\end{solution}

\begin{solution}{iq:dv:build-a-pricing-library:6}
Figer les valeurs de référence issues de la bibliothèque Python pour un ensemble de cas~; porter la même arithmétique, opération par opération (grilles, valeurs aux bords, ordre de résolution)~; tester
à $10^{-9}$ près dans les deux langages en intégration continue~; n'optimiser qu'ensuite, en conservant les tests.
\iqlookfor{valeurs de référence et équivalence opération par opération.}
\end{solution}
